\chapter*{General Conclusion}
\addcontentsline{toc}{chapter}{General Conclusion}

This project successfully addressed the complex challenge of Arabic identity matching by designing, implementing, and evaluating a modern, high-performance system. The initial problem, rooted in the linguistic nuances of Arabic names and the inefficiencies of manual verification, demanded an innovative solution that went beyond traditional methods. By integrating a multi-stage matching pipeline—combining rule-based text normalization, a custom-designed phonetic algorithm (Aramix Soundex), and a weighted scoring model with Jaro-Winkler and Levenshtein distances—we developed a system capable of delivering highly accurate results in real-time.

The architectural choices, centered on a high-performance Rust backend, a responsive Angular frontend, and a scalable PostgreSQL database, proved to be robust and effective. The system's capabilities were further extended through the integration of a Neo4j graph database for modeling complex family relationships and an n8n automation workflow that leverages generative AI for the dynamic visualization of family trees. This microservices-oriented approach not only ensured a clear separation of concerns but also provided a scalable and maintainable foundation for future development.

The project was executed following a hybrid methodology, combining the structured lifecycle of CRISP-DM for the data science components with the agile framework of Scrum for overall project management. This allowed for iterative development, continuous feedback, and the systematic alignment of technical work with business objectives. The evaluation results confirmed the system's success, demonstrating both high precision in matching and low latency under load, thereby meeting all primary success criteria.

Ultimately, this work represents more than just a technical achievement. It provides a practical blueprint for modernizing critical government infrastructure, replacing an error-prone manual process with an intelligent, efficient, and reliable solution. The system not only empowers government agents but also serves the public by ensuring that citizens can be identified quickly and accurately, restoring access and dignity. Future work could focus on integrating self-hosted large language models to enhance security and speed, expanding the training dataset, and further refining the scoring weights based on operational feedback. However, as it stands, the Tunisian_NamesML system is a complete and successful solution, ready for deployment.